\documentclass[a4paper,11pt]{article}
\usepackage[utf8]{inputenc}
\usepackage[T1]{fontenc}
\usepackage[french]{babel}
\usepackage{amsmath,amsfonts,amssymb}
\usepackage{geometry}
\usepackage{fancyhdr}
\usepackage{titlesec}
\usepackage{enumitem}

\geometry{margin=2.5cm}

% Header and footer configuration
\pagestyle{fancy}
\fancyhf{}
\fancyhead[L]{Licence 1 Informatique}
\fancyhead[R]{CERI - Avignon Université}
\fancyfoot[C]{Mathématiques Discrètes \hfill TD7 - Permutations et Combinaisons \hfill Page \thepage/3}

% Title formatting
\titleformat{\section}
{\Large\bfseries}
{\thesection}
{1em}
{}

\titleformat{\subsection}
{\large\bfseries}
{\thesubsection}
{1em}
{}

\begin{document}

\begin{center}
{\Huge\bfseries TD7 - Permutations et Combinaisons} \\[0.5cm]
{\large Hiver 2026} \\[0.3cm]
\end{center}

\vspace{0.5cm}

L'objectif de ce TD est de manipuler les permutations et combinaisons.

\section{Quelques calculs}

\begin{enumerate}
\item Combien de façon existe-t-il de sélectionner 5 éléments dans l'ordre dans un ensemble de 3 éléments si les répétitions d'éléments sont autorisées ?

\item Combien de façon existe-t-il de sélectionner 5 éléments dans l'ordre dans un ensemble de 5 éléments si les répétitions d'éléments ne sont pas autorisées ?

\item Combien de chaînes de 6 caractères peut-on former si
\begin{enumerate}
\item seules les lettres en minuscules sont autorisées ?
\item les lettres en minuscules et les lettres en majuscules sont autorisées ?
\item les lettres en minuscules, les lettres en majuscules, et les chiffres sont autorisés ?
\end{enumerate}

\item Combien de façons existe-t-il de distribuer 3 tâches à réaliser à 5 employés, si chaque employé peut effectuer 0, 1, 2 ou 3 tâches ?

\item Chaque jour, je prends un sandwich au hasard dans une pile contenant de nombreux exemplaires de 6 sandwiches différents. Si l'ordre de sélection des sandwiches compte, combien de façons existe-t-il de choisir les sandwiches pendant une semaine de 7 jours ?
\end{enumerate}

\section{À la boulangerie}

Une boulangerie propose plusieurs sortes de petits pains spéciaux. Elle vend des assortiments dégustation, où les clients choisissent le nombre de petits pains et peuvent demander certaines variétés de petits pains en particulier. La boulangerie propose les variétés suivantes : pain complet, pain aux noix, pain brioché, pain aux noisettes et raisins, pain aux aigles, pain multi-graines, pain aux figues.

\begin{enumerate}
\item Combien existe-t-il d'assortiments de 6 petits pains ?

\item Combien existe-t-il d'assortiments de 12 petits pains ?

\item Combien existe-t-il d'assortiments de 24 petits pains ?

\item Combien existe-t-il d'assortiments de 12 petits pains contenant au moins un de chaque variété ?

\item Combien existe-t-il d'assortiments de 12 petits pains contenant au moins 3 petits pains aux noix et pas plus de 2 petits pains brioché ?

\item Combien existe-t-il d'assortiments de 12 petits pains contenant aucun petit pain au noisettes et raisins ni de petits pains aux figues et contenant au moins deux petits pains de chaque autre variété ?
\end{enumerate}

\section{À la pizzeria}

Un traiteur propose de réaliser des livraisons de pizzas pour des fêtes. Ils ont à la carte les pizzas suivantes : margherita ; jambon champignons ; trois fromages ; poulet crème ; saumon pommes de terres ; aubergines.

\begin{enumerate}
\item Combien existe-t-il de commandes possible de 12 pizzas ?

\item Combien existe-t-il de commandes possible de 36 pizzas ?

\item Combien existe-t-il de commandes possible de 24 pizzas avec au moins deux de chaque ?

\item Combien existe-t-il de commandes possible de 24 pizzas contenant pas plus de deux pizzas margherita ?

\item Combien existe-t-il de commandes possible de 24 pizzas contenant au moins 5 pizzas aubergine, exactement 3 pizzas trois fromages ?

\item Combien existe-t-il de commandes possible de 24 pizzas contenant aucune pizza margherita, au moins une pizza poulet crème, au moins une pizza jambon, et aucune pizza jambon aubergine et au moins 3 pizzas trois fromages ?

\item Combien existe-t-il de commandes possible de 12 pizzas contenant au moins une pizza aubergine, au moins deux pizzas jambon champignons, au moins trois pizzas trois fromages, au moins une pizza margherita, au moins deux pizzas saumon pommes de terres ?
\end{enumerate}

\section{Anagrammes}

\begin{enumerate}
\item Combien de chaînes de caractères différentes peut-on former avec les lettres du mot \emph{JOUJOU} ?

\item Combien de chaînes de caractères différentes peut-on former avec les lettres du mot \emph{MISSISSIPPI} ?

\item Combien de chaînes de caractères différentes peut-on former avec les lettres du mot \emph{ABRACADABRANTE} ?

\item Combien de chaînes de caractères différentes peut-on former avec les lettres du mot \emph{TROTTINETTE} ?

\item Combien de chaînes de caractères différentes peut-on former avec les lettres du mot \emph{ORNITHORYNQUE} si les lettres $QU$ doivent rester ensemble dans cet ordre ?

\item Combien de chaînes de caractères différentes peut-on former avec l'ensemble des lettres de l'alphabet dans cet ordre ?
\end{enumerate}

\section{Jeu de cartes}

\begin{enumerate}
\item Cinq personnes jouent à un jeu de cartes avec un paquet standard de 52 cartes. Combien existe-t-il de façons de distribuer 7 cartes à chacun ?

\item Quatre personnes jouent à un jeu de cartes avec un paquet standard de 52 cartes. Combien existe-t-il de façons de distribuer la totalité du paquet de façon à ce que chaque joueur ait le même nombre de carte que les autres ?

\item Six personnes jouent à un jeu de cartes avec un paquet spécial de 48 cartes (toutes différentes). Combien existe-t-il de façons de distribuer 5 cartes à chacun ?
\end{enumerate}

\section{Rangement à la médiathèque}

Une médiathèque range ses rayons et s'interroge sur les différentes façons de faire.

\begin{enumerate}
\item Dans l'allée des BD, on range 12 BD sur 4 étagères différentes.
\begin{enumerate}
\item Combien de façons existe-t-il de ranger ces BD s'il s'agit de 12 exemplaires de la même BD ?
\item Combien de façons existe-t-il de ranger ces BD si ce sont 12 BD différentes et que l'ordre dans lequel elles sont agencées sur les étagères compte ?
\end{enumerate}

\item Dans l'allée des romans, on range 40 romans sur 6 étagères différentes.
\begin{enumerate}
\item Combien de façons existe-t-il de ranger ces romans s'il s'agit de 40 exemplaires du même roman ?
\item Combien de façons existe-t-il de ranger ces romans si ce sont 40 romans différents et que l'ordre dans lequel ils sont agencés sur les étagères compte ?
\end{enumerate}

\item Dans l'allée des DVD, on range 25 DVD sur 5 étagères différentes.
\begin{enumerate}
\item Combien de façons existe-t-il de ranger ces DVD s'il s'agit de 25 exemplaires du même DVD ?
\item Combien de façons existe-t-il de ranger ces DVD si ce sont 25 DVD différents et que l'ordre dans lequel ils sont agencés sur les étagères compte ?
\end{enumerate}

\item De façon générale, on range $n$ éléments sur $k$ étagères différentes.
\begin{enumerate}
\item Combien de façons existe-t-il de ranger ces éléments s'il s'agit de $n$ exemplaires du même élément ?
\item Combien de façons existe-t-il de ranger ces éléments si ce sont $n$ éléments différents et que l'ordre dans lequel ils sont agencés sur les étagères compte ?
\end{enumerate}
\end{enumerate}

\section{Manipuler des permutations et combinaisons}

\subsection{Permutations}

\begin{enumerate}
\item Donnez la prochaine permutation après 1, 4, 3, 2 dans l'ordre lexicographique

\item Donnez la prochaine permutation après 5, 4, 1, 2, 3 dans l'ordre lexicographique

\item Donnez la prochaine permutation après 1, 2, 4, 5, 3 dans l'ordre lexicographique

\item Donnez la prochaine permutation après 4, 5, 2, 3, 1 dans l'ordre lexicographique

\item Donnez la prochaine permutation après 5, 4, 3, 2, 1 dans l'ordre lexicographique

\item Donnez la prochaine permutation après 6, 7, 1, 4, 2, 3, 5 dans l'ordre lexicographique

\item Donnez la prochaine permutation après 3, 1, 5, 2, 8, 7, 6, 4 dans l'ordre lexicographique
\end{enumerate}

\subsection{Combinaisons}

\begin{enumerate}
\item Donnez toutes les 3-combinaisons de $\{1, 2, 3, 4, 5\}$ dans l'ordre lexicographique

\item Donnez toutes les 4-combinaisons de $\{1, 2, 3, 4, 5, 6\}$ dans l'ordre lexicographique
\end{enumerate}

\section{Pour aller plus loin}

\begin{enumerate}
\item Proposez un algorithme permettant d'énumérer toutes les combinaisons d'un ensemble donné.

\item Proposez un algorithme récursif pour énumérer toutes les permutations d'un ensemble donné.
\end{enumerate}

\end{document}
